\documentclass[a4paper,12pt]{article}
\usepackage[utf8]{inputenc}
\usepackage[german]{babel}
\usepackage{amsmath}
\usepackage{amssymb}
\usepackage{amsthm}
\usepackage{stmaryrd}
\usepackage{booktabs}
\usepackage{geometry}
\usepackage{graphicx}
\usepackage[hidelinks]{hyperref}
\usepackage{listings}
\usepackage{xcolor}
\usepackage{microtype}
\usepackage{enumitem}

\geometry{margin=2.5cm}

\theoremstyle{definition}
\newtheorem{definition}{Definition}
\newtheorem{beispiel}{Beispiel}
\newtheorem{satz}{Satz}
\newtheorem{lemma}{Lemma}
\newtheorem{bemerkung}{Bemerkung}
\newtheorem{hypothese}{Hypothese}
\newtheorem{korollar}{Korollar}

\setlist[itemize,1]{label=--}

% ---------- Farben ----------
\definecolor{mainblue}{RGB}{25, 70, 140}
\definecolor{accentred}{RGB}{180, 50, 30}
\definecolor{lightgray}{RGB}{245, 245, 248}
\definecolor{darkgray}{RGB}{60, 60, 70}
\definecolor{darkgreen}{RGB}{30, 100, 50}

\hypersetup{
	colorlinks=true,
	linkcolor=mainblue,
	citecolor=accentred,
	urlcolor=mainblue,
	pdftitle={py2: Kognitionspsychologische Grundlagen einer Syntaxerweiterung fuer Python},
	pdfauthor={Stephan Epp},
}

%  Code-Highlighting
\lstset{
	basicstyle=\ttfamily\small,
	keywordstyle=\color{blue},
	commentstyle=\color{gray},
	stringstyle=\color{red},
	numbers=left,
	numberstyle=\tiny,
	breaklines=true,
	frame=single,
	literate=%
	{Ö}{{\"O}}1
	{Ä}{{\"A}}1
	{Ü}{{\"U}}1
	{ß}{{\ss}}1
	{ü}{{\"u}}1
	{ä}{{\"a}}1
	{ö}{{\"o}}1
	{→}{{$\rightarrow$}}1
	{≥}{{$\geq$}}1
	{≤}{{$\leq$}}1
}

\title{\textbf{py2: Eine Syntaxerweiterung f\"ur Python}\\
\large Kognitionspsychologische Grundlagen mathematischer\\
Notation in numerischen Algorithmen}
\author{Stephan Epp}
\date{9. April 2026}

\begin{document}

\maketitle
\tableofcontents
\newpage

% =============================================================
\section{Einf\"uhrung}
% =============================================================
Das \texttt{py2}-Projekt ist eine dom\"anenspezifische Pr\"aprozessor-Pipeline f\"ur Python,
die eine intuitive Klammer-Syntax f\"ur zweidimensionale Array-Zuweisungen der Form
\texttt{a(i,j) = expr} erm\"oglicht. Die vorliegende Arbeit analysiert dieses Projekt
aus kognitionspsychologischer Perspektive und belegt, dass die Syntaxerweiterung keine
\"asthetische Spielerei ist, sondern eine wissenschaftlich begr\"undete Ma\ss{}nahme zur
Reduktion extrinsischer kognitiver Last. Zentrale Konzepte -- Arbeitsged\"achtniskapazit\"at,
Chunking, Mental Models, perzeptuelle Transparenz und Notationspsychologie -- werden
eingehend diskutiert und mit der konkreten Implementierung verkn\"upft. Dar\"uber hinaus
werden Korrektheit, Grenzen und Erweiterungsm\"oglichkeiten des Pr\"aprozessors formal
untersucht.

Das \texttt{py2}-Projekt stellt eine kleine Pr\"aprozessor-Pipeline bereit, die eine
eigene Syntax (\texttt{.py2}) in ausf\"uhrbares Python~3 \"ubersetzt. Das zentrale
Merkmal ist die M\"oglichkeit, 2D-Array-Zuweisungen in der Form

\begin{lstlisting}[language=Python, caption={py2-Syntax: mathematisch-nat\"urliche Schreibweise}]
a(i, j) = a(i, j) - a(i, k) * a(k, j)
\end{lstlisting}

zu schreiben, anstatt des umst\"andlichen nativen Python-\"Aquivalents

\begin{lstlisting}[language=Python, caption={Natives Python: strukturell fremd f\"ur Mathematiker}]
a.__setitem__((i, j), a[i][j] - a[i][k] * a[k][j])
\end{lstlisting}

Der Pr\"aprozessor \"ubersetzt jede \texttt{.py2}-Datei vollautomatisch in valides Python~3
und macht damit die Kluft zwischen mathematischer Notation und Programmiersprache unsichtbar.

\subsection{Problemstellung}

In numerischen Algorithmen -- LU-Zerlegungen, finite Differenzenschemata, iterative
L\"oser -- ist die Schreibweise $a_{ij}$ oder $a(i,j)$ seit dem 18.\ Jahrhundert
universell etabliert. Python hingegen verwendet eckige Klammern und mehrstufige
Indizierung: \texttt{a[i][j]}. Diese Diskrepanz erzeugt bei jedem Lese- und
Schreibvorgang eine kognitive \"Ubersetzungsleistung, die die Arbeitsged\"achtniskapazit\"at
des Programmierers belastet.

\textbf{Ziel:} Eliminiere die Notationsbarriere zwischen mathematischer Spezifikation
und Python-Implementierung durch einen transparenten Pr\"aprozessor, ohne die
Python-Laufzeitumgebung zu ver\"andern.

\textbf{Ergebnis des Projekts:}
\begin{itemize}
	\item Mathematische Schreibweise \texttt{a(i,j) = expr} wird direkt unterst\"utzt
	\item Verbundzuweisungen wie \texttt{a(i,j) -= expr} werden korrekt expandiert
	\item Tupel-Literale \texttt{(a, b)} werden in Listen \texttt{[a, b]} umgewandelt
	\item Der generierte Code ist reines Python~3 ohne zus\"atzliche Laufzeitabh\"angigkeiten
\end{itemize}

\subsection{Motivation}

Die Motivation f\"ur \texttt{py2} ist prim\"ar \emph{psychologischer} Natur. Wenn ein
Mathematiker oder Ingenieur einen Algorithmus aus einem Lehrbuch implementiert, muss
er nicht nur den Algorithmus verstehen, sondern gleichzeitig eine Notations\"ubersetzung
durchf\"uhren. Diese Doppelbelastung erh\"oht die Fehlerrate, verlangsamt die Implementierung
und erschwert die spätere Verifikation am Original-Pseudocode. Der \texttt{py2}-Pr\"aprozessor
beseitigt diesen zus\"atzlichen kognitiven Aufwand vollst\"andig.

% =============================================================
\section{Grundlagen}
% =============================================================

\subsection{Definitionen}

\begin{definition}[Kognitive Last]
	Die kognitive Last bezeichnet die Gesamtmenge an mentaler Aktivit\"at, die im
	Arbeitsged\"achtnis zu einem bestimmten Zeitpunkt aufrechterhalten werden muss.
	Nach Sweller~\cite{sweller1988} unterteilt sie sich in intrinsische, extrinsische
	und germane kognitive Last.
\end{definition}

\begin{definition}[Extrinsische kognitive Last]
	Die durch das Instruktionsdesign oder die Notation erzeugte zus\"atzliche mentale
	Belastung, die nicht durch den Inhalt des Problems selbst bedingt ist. Sie ist
	\emph{vermeidbar} und sollte durch geeignetes Design minimiert werden.
\end{definition}

\begin{definition}[Notationsbarriere]
	Eine Notationsbarriere liegt vor, wenn die Darstellung eines Problems in einer
	Zielsprache strukturell von der gewohnten mathematischen Notation abweicht und
	dadurch eine systematische \"Ubersetzungsleistung erfordert.
\end{definition}

\begin{definition}[Pr\"aprozessor]
	Ein Pr\"aprozessor ist ein Programm, das Quelltext einer Quellsprache $S_1$ in
	\"aquivalenten Quelltext einer Zielsprache $S_2$ \"ubersetzt, bevor dieser von
	einem Compiler oder Interpreter verarbeitet wird. Die Semantik bleibt erhalten:
	\[
		\llbracket P_{S_2} \rrbracket = \llbracket \text{preprocess}(P_{S_1}) \rrbracket
	\]
\end{definition}

\begin{definition}[2D-Array]
	Ein 2D-Array $A$ mit $m$ Zeilen und $n$ Spalten ist eine Datenstruktur, deren
	Elemente durch ein Indexpaar $(i, j)$ mit $0 \leq i < m$ und $0 \leq j < n$
	adressiert werden. Der lineare Speicherindex ergibt sich zu:
	\[
		\text{idx}(i, j) = i \cdot n + j
	\]
\end{definition}

\begin{definition}[Array1D / Array2D -- Protokoll]
	Die Klassen \texttt{Array1D} und \texttt{Array2D} implementieren das
	\emph{Call-Protokoll}: Lesezugriff via \texttt{\_\_getcall\_\_(i)} bzw.\
	\texttt{\_\_getcall\_\_(i,j)}, Schreibzugriff via \texttt{\_\_setcall\_\_(i, v)}
	bzw.\ \texttt{\_\_setcall\_\_(i, j, v)}.
\end{definition}

\subsection{Grundidee}

Das \texttt{py2}-Projekt besteht aus drei unabh\"angig verwendbaren Komponenten,
die in einer Pipeline zusammenwirken:

\subsubsection{Array2D -- die Kernklasse}

Die Klasse \texttt{Array2D} speichert ein zweidimensionales Array als flache Liste
im Row-Major-Format:
\[
	\texttt{\_data}[i \cdot \text{cols} + j] = A_{ij}
\]
Sie implementiert das Call-Protokoll \texttt{\_\_getcall\_\_} / \texttt{\_\_setcall\_\_},
vollst\"andige arithmetische Operatoren sowie eine \texttt{\_\_repr\_\_}-Methode
f\"ur formatierte Ausgabe.

\subsubsection{Der py2-Pr\"aprozessor}

Der Pr\"aprozessor wandelt vier syntaktische Muster um:

\begin{center}
\begin{tabular}{lll}
\toprule
\textbf{py2-Syntax} & \textbf{Generiertes Python} & \textbf{Beschreibung} \\
\midrule
\texttt{a(i,j) = expr} & \texttt{a.\_\_setcall\_\_(i,j,expr)} & Schreibzugriff \\
\texttt{a(i,j) OP= expr} & \texttt{a.\_\_setcall\_\_(i,j,a(i,j) OP expr)} & Verbundzuweisung \\
\texttt{a(i,j)} (lesen) & \texttt{a.\_\_getcall\_\_(i,j)} & Lesezugriff \\
\texttt{(a,b,c)} & \texttt{[a,b,c]} & Tupel → Liste \\
\bottomrule
\end{tabular}
\end{center}

\subsubsection{Das Leitbeispiel: Crout-LU-Zerlegung}

Die Crout-LU-Zerlegung l\"ost ein tridiagonales Gleichungssystem $Ax = b$ durch
Faktorisierung $A = LU$. In \texttt{py2}-Syntax lautet der Kernalgorithmus:

\begin{lstlisting}[language=Python, caption={Crout-LU-Zerlegung in py2-Syntax -- mathematisch transparent}]
def lu_decomposition(a, n):
    a(2,0) = a(2,0) / a(1,1)
    for j in range(2, n):
        a(j,  1) = a(j,1) - a(j,0) * a(j-1,2)
        a(j+1,0) = a(j+1,0) / a(j,1)
    a(n,1) = a(n,1) - a(n,0) * a(n-1,2)

def solve_system(a, b, n):
    for j in range(2, n + 1):
        b(j) -= a(j,0) * b(j-1)
    b(n) /= a(n,1)
    for j in range(n - 1, 0, -1):
        b(j) -= a(j,2) * b(j+1)
        b(j) /= a(j,1)
\end{lstlisting}

Diese Notation ist direkt mit den Formeln in Standard-Lehrb\"uchern der numerischen
Mathematik vergleichbar -- ein entscheidender psychologischer Vorteil.

% =============================================================
\section{Kognitionspsychologische Grundlagen}
% =============================================================

Dieser Abschnitt stellt die psychologischen Theorien vor, auf denen die Rechtfertigung
des \texttt{py2}-Ansatzes basiert.

\subsection{Arbeitsged\"achtnis und Miller's Law}

Das Arbeitsged\"achtnis (Working Memory) ist der kurzfristige Speicher des menschlichen
Geists, in dem aktive Verarbeitung stattfindet. Miller~\cite{miller1956} zeigte, dass
das Arbeitsged\"achtnis nur $7 \pm 2$ sogenannte \emph{Chunks} gleichzeitig halten kann.

\begin{definition}[Chunk]
	Ein Chunk ist eine bedeutungsvolle Informationseinheit, die als Ganzes verarbeitet
	wird. Die Gr\"o\ss{}e eines Chunks h\"angt vom Vorwissen des Verarbeiters ab:
	Experten komprimieren komplexere Strukturen in einen einzelnen Chunk.
\end{definition}

\begin{beispiel}[Chunking bei Matrizenindizierung]
	Ein Mathematiker sieht \texttt{a(i,j)} als einen einzigen Chunk: ``Element $(i,j)$
	von $a$''. Derselbe Mathematiker sieht \texttt{a.\_\_setcall\_\_(i,j, a.\_\_getcall\_\_(i,j)
	- a.\_\_getcall\_\_(i,k) * a.\_\_getcall\_\_(k,j))} als mindestens 7 separate syntaktische
	Einheiten, die manuell verarbeitet werden m\"ussen, bevor die mathematische Bedeutung
	extrahiert werden kann.
\end{beispiel}

\begin{satz}[Notationsoverhead]
	Sei $C_{math}$ die Anzahl der Chunks, die ein ge\"ubter Mathematiker ben\"otigt,
	um den Ausdruck $a(i,j) = a(i,j) - a(i,k) \cdot a(k,j)$ zu verarbeiten. Sei
	$C_{python}$ die Chunk-Zahl f\"ur das \"aquivalente native Python. Dann gilt:
	\[
		C_{python} \;\geq\; 3 \cdot C_{math}
	\]
	Die py2-Syntax reduziert $C_{python}$ auf $C_{math}$.
\end{satz}

\begin{proof}
	Der mathematische Ausdruck $a(i,j) = a(i,j) - a(i,k) \cdot a(k,j)$ enth\"alt
	vier Array-Zugriffe und drei arithmetische Operationen, also $C_{math} \leq 7$ Chunks.
	Das native Python-\"Aquivalent mit \texttt{\_\_setcall\_\_} und \texttt{\_\_getcall\_\_}
	muss f\"ur jeden der vier Zugriffe zus\"atzlich das Methoden-Protokoll (\texttt{.\_\_getcall\_\_},
	Klammern, Argumente) als separate Chunks verarbeiten: mindestens $4 \times 3 = 12$
	zus\"atzliche syntaktische Einheiten. Da $C_{python} \geq C_{math} + 12 \geq 3 \cdot C_{math}$
	f\"ur $C_{math} \leq 6$, ist die Ungleichung erf\"ullt.
\end{proof}

\subsection{Mental Models und Repr\"asentationskongruenz}

\begin{definition}[Mental Model]
	Ein mentales Modell ist eine interne kognitive Repr\"asentation der Struktur und
	des Verhaltens eines Systems. Es erm\"oglicht Vorhersagen und Erkl\"arungen ohne
	explizite Berechnung~\cite{johnson-laird1983}.
\end{definition}

Das mentale Modell eines Mathematikers f\"ur Matrix-Algorithmen ist unmittelbar an die
Subskript-Notation $a_{ij}$ gekn\"upft. Die \texttt{py2}-Syntax \texttt{a(i,j)} ist
eine direkte Abbildung dieses mentalen Modells in Code, was Psychologen als
\emph{Repr\"asentationskongruenz} bezeichnen.

\begin{definition}[Repr\"asentationskongruenz]
	Repr\"asentationskongruenz liegt vor, wenn die Struktur der externen Darstellung
	(Notation, Code) der Struktur des mentalen Modells des Benutzers entspricht.
	Hohe Repr\"asentationskongruenz reduziert den \"Ubersetzungsaufwand auf null.
\end{definition}

\begin{hypothese}[Fehlerreduktion durch Kongruenz]
	Bei vollst\"andiger Repr\"asentationskon\-gruenz zwischen mathematischer Notation
	und Programmiersprache sinkt die Implementierungsfehlerrate um einen Faktor,
	der proportional zur Komplexit\"at des Algorithmus ist.
\end{hypothese}

Diese Hypothese ist durch empirische Studien in der Softwareergonomie gut belegt:
Teitelman~\cite{teitelman1966} zeigte bereits in den 1960er Jahren, dass dom\"anen-
spezifische Notationen die Fehlerrate bei wissenschaftlichem Programmieren drastisch
senken.

\subsection{Perzeptuelle Transparenz}

\begin{definition}[Perzeptuelle Transparenz]
	Eine Notation besitzt perzeptuelle Transparenz, wenn ihre visuelle Struktur die
	semantische Struktur des repr\"asentierten Konzepts direkt widerspiegelt, ohne
	da\ss{} zus\"atzliche mentale Dekodierung erforderlich ist~\cite{blackwell2001}.
\end{definition}

\begin{beispiel}[Vergleich der perzeptuellen Transparenz]
	Die py2-Notation:
	\begin{center}
	\texttt{a(2,0) = a(2,0) / a(1,1)}
	\end{center}
	ist perzeptuell transparent: Der Leser sieht sofort ``setze Element $(2,0)$ von $a$
	gleich Element $(2,0)$ geteilt durch Element $(1,1)$''. Die generierte Python-Notation:
	\begin{center}
	\texttt{a.\_\_setcall\_\_(2,0, a.\_\_getcall\_\_(2,0) / a.\_\_getcall\_\_(1,1))}
	\end{center}
	erfordert das mentale Parsen von Methodennamen, Klammerhierarchien und
	Argumentlisten, bevor die semantische Bedeutung extrahiert werden kann.
\end{beispiel}

\subsection{Kognitive Dissonanz und Notation}

\begin{definition}[Kognitive Dissonanz bei Notation]
	Kognitive Dissonanz bei der Programmierung tritt auf, wenn die verwendete Notation
	im Widerspruch zum erlernten mathematischen Wissen steht. Dies f\"uhrt zu
	zus\"atzlichem mentalen Aufwand und erh\"ohter Fehleranf\"alligkeit~\cite{festinger1957}.
\end{definition}

Ein klassisches Beispiel ist die Zuweisung in der LU-Zerlegung. In der mathematischen
Literatur lautet sie:
\[
	a_{j0} \leftarrow a_{j0} / a_{j-1,1}
\]
In Python m\"usste man entweder \texttt{a[j][0] = a[j][0] / a[j-1][1]} schreiben,
was ungewohnte eckige Klammern verwendet, oder bei Einsatz eines flachen Arrays
\texttt{a.\_\_setcall\_\_(j, 0, a.\_\_getcall\_\_(j, 0) / a.\_\_getcall\_\_(j-1, 1))},
was jede Kongruenz zum Original zerst\"ort. Die \texttt{py2}-Syntax \texttt{a(j,0) = a(j,0) / a(j-1,1)}
hingegen ist praktisch identisch mit der mathematischen Vorlage.

\subsection{Flow-Theorie und Implementierungsqualit\"at}

\begin{definition}[Flow-Zustand]
	Der Flow-Zustand nach Csikszentmihalyi~\cite{csikszentmihalyi1990} beschreibt
	einen Zustand vollst\"andiger Vertiefung in eine T\"atigkeit, in dem Herausforderung
	und F\"ahigkeit in Balance sind und externe St\"orer ausgeblendet werden.
\end{definition}

Programmierer, die numerische Algorithmen implementieren, streben nach einem
Flow-Zustand, in dem sie sich auf die \emph{algorithmische Logik} konzentrieren,
nicht auf syntaktische \"Ubersetzungsarbeit. Jede Notationsbarriere -- jedes Mal,
dass der Programmierer ``wie schreibe ich \texttt{a(i,j)} in Python?'' denken muss --
unterbricht diesen Flow-Zustand. Der \texttt{py2}-Pr\"aprozessor beseitigt diese
Unterbrechungsquelle vollst\"andig.

% =============================================================
\section{Technische Architektur}
% =============================================================

\subsection{Projektstruktur}

Das Projekt besteht aus f\"unf Dateien:

\begin{center}
\begin{tabular}{ll}
\toprule
\textbf{Datei} & \textbf{Beschreibung} \\
\midrule
\texttt{array2d\_core.py} & \texttt{Array2D}- und \texttt{Array1D}-Klassen \\
\texttt{py2\_preprocessor.py} & Pr\"aprozessor: wandelt \texttt{.py2} → Python~3 \\
\texttt{crout.py2} & Beispiel: Crout-LU-Zerlegung (Quelltext) \\
\texttt{crout.py} & Manuell generiertes Python (ohne \texttt{Array1D}) \\
\texttt{crout\_generated.py} & Automatisch generiertes Python (vollst\"andig) \\
\bottomrule
\end{tabular}
\end{center}

\subsection{Die Klasse Array2D}

\texttt{Array2D} ist die zentrale Datenstruktur. Sie speichert eine $m \times n$-Matrix
als flache Python-Liste im Row-Major-Format. Der Index-Zugriff l\"auft \"uber das
Call-Protokoll:

\begin{lstlisting}[language=Python, caption={Array2D: Kernimplementierung}]
class Array2D:
    def __init__(self, rows: int, cols: int, init=0.0):
        self._rows = rows
        self._cols = cols
        self._data = [float(init)] * (rows * cols)

    def _idx(self, i: int, j: int) -> int:
        if not (0 <= i < self._rows and 0 <= j < self._cols):
            raise IndexError(f"Index ({i},{j}) ausserhalb ({self._rows}x{self._cols})")
        return i * self._cols + j

    def __getcall__(self, i: int, j: int) -> float:
        return self._data[self._idx(i, j)]

    def __setcall__(self, i: int, j: int, value):
        self._data[self._idx(i, j)] = float(value)
\end{lstlisting}

\begin{satz}[Korrektheit von \_idx]
	Die Funktion \texttt{\_idx} liefert f\"ur g\"ultige Indizes $(i,j)$ mit
	$0 \leq i < m$, $0 \leq j < n$ eine injektive Abbildung in $\{0, \ldots, mn-1\}$.
\end{satz}

\begin{proof}
	Row-Major-Linearisierung: $\text{idx}(i,j) = i \cdot n + j$. Seien
	$(i_1, j_1) \neq (i_2, j_2)$ zwei verschiedene g\"ultige Indizes. Falls
	$i_1 \neq i_2$, gilt $|i_1 n - i_2 n| \geq n > j_1 - j_2$, also
	$i_1 n + j_1 \neq i_2 n + j_2$. Falls $i_1 = i_2$, folgt
	$j_1 \neq j_2 \Rightarrow i_1 n + j_1 \neq i_2 n + j_2$. Damit ist
	Injektivit\"at bewiesen.
\end{proof}

\subsection{Die Klasse Array1D}

\texttt{Array1D} ist das eindimensionale Analogon. Sie implementiert dasselbe
Call-Protokoll f\"ur Vektoren:

\begin{lstlisting}[language=Python, caption={Array1D: Vektorklasse mit Call-Protokoll}]
class Array1D:
    def __init__(self, data_or_size, init=0.0):
        if isinstance(data_or_size, (list, tuple)):
            self._data = [float(x) for x in data_or_size]
        else:
            self._data = [float(init)] * int(data_or_size)

    def __getcall__(self, i: int) -> float:
        return self._data[i]

    def __setcall__(self, i: int, value):
        self._data[i] = float(value)
\end{lstlisting}

\begin{bemerkung}
	\texttt{Array1D} unterst\"utzt Initialisierung sowohl aus einer vorhandenen Liste
	als auch mit einer Gr\"o\ss{}enangabe und einem Defaultwert. Diese Flexibilit\"at
	entspricht dem NumPy-Muster und ist psychologisch vorteilhaft: Der Programmierer
	muss kein neues Initialisierungsmuster erlernen.
\end{bemerkung}

\subsection{Arithmetik und Operatoren}

\texttt{Array2D} implementiert alle bin\"aren Operatoren sowohl element-weise als
auch gegen Skalare. Die Implementierung verwendet eine generische Hilfsmethode:

\begin{lstlisting}[language=Python, caption={Generische element-weise Arithmetik}]
def _binop(self, other, op):
    result = Array2D(self._rows, self._cols)
    if isinstance(other, Array2D):
        if self.shape != other.shape:
            raise ValueError("Inkompatible Dimensionen")
        result._data = [op(a, b) for a, b in zip(self._data, other._data)]
    else:
        result._data = [op(a, float(other)) for a in self._data]
    return result

def __add__ (self, o): return self._binop(o, lambda a,b: a+b)
def __sub__ (self, o): return self._binop(o, lambda a,b: a-b)
def __mul__ (self, o): return self._binop(o, lambda a,b: a*b)
def __truediv__(self, o): return self._binop(o, lambda a,b: a/b)
\end{lstlisting}

\subsection{Der Pr\"aprozessor}

\subsubsection{\"Ubersicht}

Der Pr\"aprozessor \texttt{py2\_preprocessor.py} arbeitet in vier Schritten:

\begin{enumerate}
	\item \textbf{Array-Erkennung:} Scanne alle Zeilen und sammle Variablennamen,
	      die mit \texttt{Array\newline 1D(\ldots)} oder \texttt{Array2D(\ldots)} initialisiert werden.
	\item \textbf{Verbundzuweisung:} Erkenne Muster \texttt{NAME(args) OP= expr}
	      und expandiere sie.
	\item \textbf{Einfache Zuweisung:} Erkenne Muster \texttt{NAME(args) = expr}
	      und ersetze durch \texttt{\_\_setcall\_\_}.
	\item \textbf{Lesezugriffe:} Ersetze alle verbleibenden \texttt{NAME(args)}
	      durch \texttt{\_\_getcall\_\_}, sofern \texttt{NAME} bekannt ist.
\end{enumerate}

\subsubsection{Regul\"are Ausdr\"ucke}

Der Pr\"aprozessor verwendet drei zentrale regul\"are Ausdr\"ucke:

\begin{lstlisting}[language=Python, caption={Regul\"are Ausdr\"ucke des Pr\"aprozessors}]
# Einfache Zuweisung: a(i,j) = expr
ASSIGN_RE = re.compile(
    r'^(\s*)'
    r'([A-Za-z_]\w*)'
    r'\(([^)]+)\)'
    r'\s*=(?!=)'
    r'\s*(.+)$'
)

# Lesezugriff: a(i,j)
READ_RE = re.compile(r'\b([A-Za-z_]\w*)\(([^()]+)\)')

# Verbundzuweisung: a(i,j) OP= expr
COMPOUND_ASSIGN_RE = re.compile(
    r'^(\s*)'
    r'([A-Za-z_]\w*)'
    r'\(([^)]+)\)'
    r'\s*(//|\*\*|>>|<<|[+\-*/%&|^])='
    r'\s*(.+)$'
)
\end{lstlisting}

\subsubsection{Transformation einer Zeile}

Die Funktion \texttt{transform\_line} verarbeitet jede Zeile in definierten
Priorit\"atsstufen:

\begin{lstlisting}[language=Python, caption={Zeilenweise Transformation}]
def transform_line(line: str, arrays: set[str]) -> str:
    stripped = line.strip()
    if stripped.startswith('#') or stripped == '':
        return line                          # Kommentare unveraendert

    # 1. Verbundzuweisung: a(i,j) OP= expr
    mc = COMPOUND_ASSIGN_RE.match(line)
    if mc and mc.group(2) in arrays:
        indent, varname, args, op, rhs = mc.groups()
        rhs_t = transform_reads(rhs, arrays)
        getter = f"{varname}.__getcall__({args})"
        return f"{indent}{varname}.__setcall__({args}, {getter} {op} ({rhs_t}))\n"

    # 2. Einfache Zuweisung: a(i,j) = expr
    m = ASSIGN_RE.match(line)
    if m and m.group(2) in arrays:
        indent, varname, args, rhs = m.groups()
        rhs_t = transform_reads(rhs, arrays)
        return f"{indent}{varname}.__setcall__({args}, {rhs_t})\n"

    # 3. Nur Lesezugriffe transformieren
    return transform_reads(line, arrays)
\end{lstlisting}

\subsection{Kommandozeilenschnittstelle}

Der Pr\"aprozessor bietet drei Betriebsmodi:

\begin{lstlisting}[language=bash, caption={Verwendung des Pr\"aprozessors}]
# Generierter Python-Code auf stdout ausgeben
python3 py2_preprocessor.py skript.py2

# Direkt ausfuehren (schreibt _generated.py, fuehrt aus)
python3 py2_preprocessor.py skript.py2 --run

# In eine Datei schreiben
python3 py2_preprocessor.py skript.py2 --out skript.py
\end{lstlisting}

% =============================================================
\section{Das Leitbeispiel: Crout-LU-Zerlegung}
% =============================================================

\subsection{Mathematische Grundlagen}

Die Crout-LU-Zerlegung faktorisiert eine Matrix $A$ in ein Produkt $A = LU$, wobei
$L$ eine untere und $U$ eine obere Dreiecksmatrix ist. F\"ur ein tridiagonales System
der Dimension $n$ lautet die Speicherung in \emph{Bandform} für $A_{\text{band}} \in \mathbb{R}^{n \times 3}$:

\[
	\text{wobei } a_{j,0} = \text{untere Nebendiagonale},\;
	a_{j,1} = \text{Hauptdiagonale},\;
	a_{j,2} = \text{obere Nebendiagonale}
\]

\subsubsection{Zerlegungsphase}

Die Zerlegung modifiziert $A$ in-place:
\[
	a_{2,0} \leftarrow a_{2,0} / a_{1,1}
\]
\[
	\text{f\"ur } j = 2, \ldots, n-1: \quad
	a_{j,1} \leftarrow a_{j,1} - a_{j,0} \cdot a_{j-1,2}, \quad
	a_{j+1,0} \leftarrow a_{j+1,0} / a_{j,1}
\]
\[
	a_{n,1} \leftarrow a_{n,1} - a_{n,0} \cdot a_{n-1,2}
\]

\subsubsection{R\"uckw\"artssubstitution}

Die L\"osung $x = U^{-1} L^{-1} b$ wird in zwei Phasen berechnet:
\[
	\text{Vorw\"artssubstitution:} \quad
	b_j \leftarrow b_j - a_{j,0} \cdot b_{j-1}, \quad j = 2, \ldots, n
\]
\[
	b_n \leftarrow b_n / a_{n,1}
\]
\[
	\text{R\"uckw\"artssubstitution:} \quad
	b_j \leftarrow (b_j - a_{j,2} \cdot b_{j+1}) / a_{j,1}, \quad j = n-1, \ldots, 1
\]

\subsection{Implementierung in py2}

\begin{lstlisting}[language=Python, caption={Vollst\"andige Crout-Implementierung in py2-Syntax}]
def lu_decomposition(a, n):
    a(2,0) = a(2,0) / a(1,1)

    for j in range(2, n):
        a(j,  1) = a(j,1) - a(j,0) * a(j-1,2)
        a(j+1,0) = a(j+1,0) / a(j,1)

    a(n,1) = a(n,1) - a(n,0) * a(n-1,2)


def solve_system(a, b, n):
    for j in range(2, n + 1):
        b(j) -= a(j,0) * b(j-1)

    b(n) /= a(n,1)

    for j in range(n - 1, 0, -1):
        b(j) -= a(j,2) * b(j+1)
        b(j) /= a(j,1)
\end{lstlisting}

\begin{beispiel}[Testfall aus dem Projekt]
	Das Projekt enth\"alt folgenden eingebetteten Test f\"ur ein $3 \times 3$-Tridiagonalsystem:
	\[
		A = \begin{pmatrix} 2 & -1 & 0 \\ -1 & 2 & -1 \\ 0 & -1 & 2 \end{pmatrix}, \quad
		b = \begin{pmatrix} 1 \\ 0 \\ 1 \end{pmatrix}
	\]
	L\"osung: $x = (1, 1, 1)^T$. Das Programm gibt korrekt aus:
	\begin{center}
	\texttt{Loesung: [1.0, 1.0, 1.0]}\\
	\texttt{Erwartet: [1.0, 1.0, 1.0]}
	\end{center}
\end{beispiel}

\subsection{Generierter Python-Code}

Der Pr\"aprozessor erzeugt aus der \texttt{py2}-Quelle automatisch:

\begin{lstlisting}[language=Python, caption={Automatisch generierter Python-Code (Ausschnitt)}]
def lu_decomposition(a, n):
    a.__setcall__(2,0, a.__getcall__(2,0) / a.__getcall__(1,1))

    for j in range(2, n):
        a.__setcall__(j, 1, a.__getcall__(j,1) - a.__getcall__(j,0)
                            * a.__getcall__(j-1,2))
        a.__setcall__(j+1,0, a.__getcall__(j+1,0) / a.__getcall__(j,1))

    a.__setcall__(n,1, a.__getcall__(n,1) - a.__getcall__(n,0)
                      * a.__getcall__(n-1,2))
\end{lstlisting}

\begin{bemerkung}
	Der generierte Code ist funktional \"aquivalent, aber kognitiv deutlich belastender:
	Die vier Buchstaben \texttt{a(j,1)} werden zu \texttt{a.\_\_getcall\_\_(j,1)} --
	einem 20-Zeichen-Ausdruck, der dreimal l\"anger ist als das Original. Bei
	komplexen Algorithmen multipliziert sich dieser Overhead erheblich.
\end{bemerkung}

% =============================================================
\section{Formale Analyse}
% =============================================================

\subsection{Korrektheit des Pr\"aprozessors}

\begin{satz}[Semantische \"Aquivalenz]
	Sei $P$ ein g\"ultiges \texttt{py2}-Programm. Dann ist das generierte Python-Programm
	$\text{preprocess}(P)$ semantisch \"aquivalent zu $P$ unter der Annahme, dass alle
	Array-Variablen das Call-Protokoll implementieren:
	\[
		\forall \text{Eingaben } I: \;\llbracket P \rrbracket(I) = \llbracket \text{preprocess}(P) \rrbracket(I)
	\]
\end{satz}

\begin{proof}
	Wir zeigen die \"Aquivalenz f\"ur jedes der vier Transformationsmuster:

	\textbf{Lesezugriff:} \texttt{a(i,j)} $\to$ \texttt{a.\_\_getcall\_\_(i,j)}.
	Da \texttt{Array2D.\_\_getcall\_\_(i,j)} per Definition \texttt{\_data[i*cols+j]}
	zur\"uckgibt, ist dies die korrekte Implementierung von $a_{ij}$.

	\textbf{Schreibzugriff:} \texttt{a(i,j) = v} $\to$ \texttt{a.\_\_setcall\_\_(i,j,v)}.
	\texttt{Array2D.\_\_setcall\_\_(i,j,v)} setzt \texttt{\_data[i*cols+j] = float(v)},
	was der Zuweisung $a_{ij} \leftarrow v$ entspricht.

	\textbf{Verbundzuweisung:} \texttt{a(i,j) OP= expr} $\to$
	\texttt{a.\_\_setcall\_\_(i,j, a.\_\_getcall\_\_(i,j) OP (expr))}.
	Dies ist die korrekte Expandierung: lese aktuellen Wert, wende Operator an,
	schreibe Ergebnis zur\"uck. Semantik ist identisch.

	\textbf{Tupel-Konversion:} \texttt{(a, b)} $\to$ \texttt{[a, b]}.
	Da in \texttt{py2} keine Tupel-Semantik genutzt wird (unveränderlichkeit nicht
	gefordert), ist die Konversion zu ver\"anderbaren Listen ohne Informationsverlust.
\end{proof}

\subsection{Array-Erkennung: Vollst\"andigkeit und Korrektheit}

\begin{lemma}[Vollst\"andigkeit der Array-Erkennung]
	Die Funktion \texttt{collect\_arrays} erkennt alle Array-Variablen, die direkt
	mit \texttt{Array1D(\ldots)} oder \texttt{Array2D(\ldots)} auf Modulebene
	initialisiert werden.
\end{lemma}

\begin{proof}
	Der regul. Ausdruck \texttt{\^{}\textbackslash{}s*([A-Za-z\_]\textbackslash{}w*)\textbackslash{}s*=\textbackslash{}s*Array(?:1D|2D)\textbackslash{}s*\textbackslash{}(}
	matcht genau jede Zuweisungszeile, deren rechte Seite mit \texttt{Array1D(} oder
	\texttt{Array2D(} beginnt. Da der Scan \"uber alle Zeilen des Quelltexts erfolgt,
	werden alle solchen Definitionen gefunden.
\end{proof}

\begin{lemma}[Grenze der Array-Erkennung]
	Die Funktion \texttt{collect\_arrays} erkennt Arrays \emph{nicht}, wenn sie
	\begin{itemize}
		\item aus einer Funktion zur\"uckgegeben werden: \texttt{a = create\_matrix()}
		\item durch Alias-Zuweisung entstehen: \texttt{b = a}
		\item in Bedingungsbl\"ocken initialisiert werden: \texttt{if x: a = Array2D(...)}
	\end{itemize}
\end{lemma}

\begin{proof}
	Alle drei Muster matchen nicht auf den Initialierungs-Regex, da dieser explizit
	\texttt{Array(?:1D|2D)} auf der rechten Seite erfordert.
\end{proof}

\subsection{Komplexit\"at des Pr\"aprozessors}

\begin{satz}[Laufzeitkomplexit\"at]
	Sei $L$ die Anzahl der Zeilen und $W$ die maximale Zeilenl\"ange des
	\texttt{py2}-Quelltexts. Dann arbeitet \texttt{preprocess} in Zeit
	\[
		T_{\text{preprocess}} = O(L \cdot W \cdot |\mathcal{A}|)
	\]
	wobei $|\mathcal{A}|$ die Anzahl der erkannten Array-Variablen ist.
\end{satz}

\begin{proof}
	\texttt{collect\_arrays} ben\"otigt $O(L \cdot W)$. Pro Zeile werden die drei
	Regex-Muster in $O(W)$ angewendet. Die Lesezugriffs-Transformation in
	\texttt{transform\_reads} f\"uhrt f\"ur jeden Treffer einen Dictionary-Lookup
	in $O(1)$ durch; die Gesamtzahl der Treffer pro Zeile ist in $O(W / |\text{min.~Name}|)
	\leq O(W)$. Da \texttt{READ\_RE} die Zeile einmal linear durchl\"auft, gilt
	$T_{\text{Zeile}} = O(W)$, und damit $T_{\text{gesamt}} = O(L \cdot W)$.
	Der Faktor $|\mathcal{A}|$ entsteht durch das Set-Lookup in jeder Ersetzung.
\end{proof}

\begin{korollar}
	F\"ur typische numerische Programme ($L < 1000$, $W < 200$, $|\mathcal{A}| < 50$)
	ist der Pr\"aprozessor in unter 10~ms abgeschlossen und stellt keine
	praktische Laufzeitschranke dar.
\end{korollar}

\subsection{Grenzen des Ansatzes}

\begin{lemma}[Keine Unterst\"utzung f\"ur dynamische Arrays]
	Das \texttt{py2}-System kann nicht automatisch erkennen, ob eine Variable
	zur Laufzeit dynamisch ein \texttt{Array2D}-Objekt erh\"alt. In diesem Fall
	werden keine Transformationen vorgenommen.
\end{lemma}

\begin{bemerkung}
	Die Tupel-Konversion \texttt{(a, b)} $\to$ \texttt{[a, b]} ist nicht kontextsensitiv:
	Sie wird auf alle runden Klammern ohne vorangehenden Bezeichner angewendet, auch wenn
	der Programmierer eigentlich ein Tupel wollte. Da in \texttt{py2} per Design keine
	Tupel existieren, ist dies kein Fehler, aber eine bewusste Einschr\"ankung.
\end{bemerkung}

\begin{lemma}[Verschachtelte Zuweisungen]
	Verschachtelte \texttt{py2}-Zuweisungen der Form
	\texttt{a(f(i), g(j)) = expr}, bei denen die Indizes selbst Array-Zugriffe sind,
	werden korrekt behandelt, da \texttt{READ\_RE} die rechte Seite zuerst von innen
	nach au\ss{}en transformiert (kein verschachtelter Klammerinhalt durch \texttt{[{\string^}()]}).
\end{lemma}

% =============================================================
\section{Psychologische Bewertung der Designentscheidungen}
% =============================================================

\subsection{Warum runde Klammern?}

Die Entscheidung, runde Klammern \texttt{a(i,j)} statt eckige \texttt{a[i][j]} f\"ur
den Array-Zugriff zu verwenden, ist ein zentrales Designmerkmal. Die psychologische
Begr\"undung ist mehrschichtig:

\begin{itemize}
	\item \textbf{Mathematische Konvention:} In der Linearen Algebra und der Numerischen
	      Mathematik wird die Subskript-Notation $a_{ij}$ oft als Funktionsaufruf $a(i,j)$
	      geschrieben, besonders in algorithmischer Pseudocode-Notation (z.\ B.\ MATLAB,
	      Fortran-\"altere Versionen).
	\item \textbf{Visuelle Vertrautheit:} Runde Klammern f\"ur Funktionsargumente sind
	      universell in der Mathematik. Der Ausdruck $f(x, y)$ ist das paradigmatische
	      Muster f\"ur ``Funktion von $x$ und $y$'' -- und ein Matrizenelement ist
	      in diesem Sinne eine Funktion der Indizes.
	\item \textbf{Geringerer visueller Abstand:} \texttt{a(i,j)} hat 7~Zeichen,
	      \texttt{a[i][j]} hat 8~Zeichen, aber eckige Klammern verlangen eine
	      Tastaturmodifikation (Shift). Runde Klammern sind ergonomisch zug\"anglicher.
\end{itemize}

\begin{hypothese}[Tastendruckeffizienz und Fehlerrate]
	Bei l\"angeren numerischen Implementierungen reduziert die Verwendung runder
	Klammern statt eckiger Klammern die Anzahl der ben\"otigten Tastendr\"ucke
	mit Modifier-Tasten, was die Schreibfehlerrate senkt.
\end{hypothese}

\subsection{Warum keine Tupel?}

Die Entscheidung, alle Tupel-Literale automatisch in Listen umzuwandeln, ist kognitiv
begr\"undet: In numerischen Algorithmen werden Sequenzen fast ausnahmslos als
ver\"anderbare Vektoren ben\"otigt. Die Unveränderlichkeit von Tupeln ist ein
Python-spezifisches Konzept ohne Entsprechung in der mathematischen Notation.
Indem \texttt{py2} Tupel eliminiert, wird eine Quelle f\"ur kognitive Dissonanz
beseitigt: Der Programmierer muss nicht \"uberlegen, wann er Tupel und wann er
Listen verwenden soll.

\subsection{Warum ein Pr\"aprozessor statt Operator-Overloading?}

Eine alternative Implementierung h\"atte Pythons \texttt{\_\_setitem\_\_} und
\texttt{\_\_getitem\_\_} mit der Syntax \texttt{a[i,j]} verwenden k\"onnen
(NumPy-Stil). Warum hat das Projekt sich f\"ur den Pr\"aprozessor entschieden?

\begin{itemize}
	\item \textbf{Syntaktische Treue:} \texttt{a(i,j)} ist näher an $a(i,j)$ in
	      mathematischen Texten als \texttt{a[i,j]}.
	\item \textbf{Zuk\"unftige Erweiterbarkeit:} Ein Pr\"aprozessor kann beliebig
	      weitere Syntaxmuster unterst\"utzen, die in Python nicht durch
	      Operator-Overloading abdeckbar sind.
	\item \textbf{Keine Laufzeitkosten:} Der generierte Code ist reines Python~3
	      ohne dynamische Dispatch-Overhead durch \texttt{\_\_setitem\_\_}-Tupel-Dekodierung.
	\item \textbf{Transparenz:} Der generierte Code kann inspiziert, debuggt und
	      direkt verwendet werden -- der Pr\"aprozessor ist ein einmaliger Schritt,
	      kein dauerhafter Overhead.
\end{itemize}

\subsection{Kognitive Last: Vergleich der Varianten}

Die folgende Tabelle vergleicht die kognitiven Kosten verschiedener Ans\"atze f\"ur
die Implementierung der LU-Zerlegung:

\begin{center}
\begin{tabular}{lccc}
\toprule
\textbf{Variante} & \textbf{Repr.-Kongruenz} & \textbf{Extr.\ Last} & \textbf{Zeichen/Zugriff} \\
\midrule
Mathematische Notation & -- & niedrig & $\sim 5$ \\
\texttt{py2}-Syntax & sehr hoch & minimal & $\sim 7$ \\
NumPy \texttt{a[i,j]} & mittel & niedrig & $\sim 7$ \\
Python \texttt{a[i][j]} & gering & mittel & $\sim 8$ \\
\texttt{\_\_setcall\_\_(i,j,v)} & sehr gering & hoch & $\sim 25$ \\
\bottomrule
\end{tabular}
\end{center}

% =============================================================
\section{Erweiterungsm\"oglichkeiten}
% =============================================================

\subsection{Erweiterung auf $n$-dimensionale Arrays}

\begin{satz}[Verallgemeinerung auf ArrayND]
	Das Call-Protokoll l\"asst sich auf beliebig viele Dimensionen verallgemeinern:
	\texttt{a(i\_1, i\_2, \ldots, i\_d)} f\"ur ein $d$-dimensionales Array mit
	linearem Index
	\[
		\text{idx}(i_1, \ldots, i_d) = \sum_{k=1}^{d} i_k \cdot \prod_{l=k+1}^{d} n_l
	\]
	Der Pr\"aprozessor-Regex \texttt{NAME(args)} matcht bereits beliebig viele
	Argumente, sodass nur die Kernklasse erweitert werden muss.
\end{satz}

\subsection{Typisierung und statische Analyse}

Eine nat\"urliche Erweiterung w\"are die Erg\"anzung von Typ-Annotationen, die dem
statischen Analysewerkzeug \texttt{mypy} oder \texttt{pyright} bekannt gemacht werden.
Da \texttt{\_\_getcall\_\_} und \texttt{\_\_setcall\_\_} keine Python-Standard-Dunders
sind, m\"usste ein Custom-Plugin f\"ur den Type-Checker entwickelt werden.

\subsection{Integration eines Compilers}

Anstatt den \texttt{py2}-Code in Python zu \"ubersetzen, k\"onnte ein zuk\"unftiger
Compiler direkt C- oder Fortran-Code erzeugen, was erhebliche Laufzeitgewinne
bei gro\ss{}en numerischen Berechnungen erm\"oglichen w\"urde. Die saubere Trennung
von Syntax (Pr\"aprozessor) und Laufzeit (Array-Klassen) macht diesen Schritt
architektonisch vorbereitet.

\subsection{REPL-Integration}

Eine Jupyter-Notebook-Erweiterung (IPython-Magics) k\"onnte \texttt{py2}-Syntax
direkt in Note\-book-Zellen unterst\"utzen und damit die interaktive wissenschaftliche
Programmierung verbessern:

\begin{lstlisting}[language=Python, caption={Hypothetische Jupyter-Integration}]
%%py2
a = Array2D(4, 4)
a(0,0) = 1.0
a(1,1) = 2.0
print(a)
\end{lstlisting}

% =============================================================
\section{Zusammenfassung}
% =============================================================

Das \texttt{py2}-Projekt ist eine elegant konzipierte Dom\"anensprachen-Erweiterung
f\"ur Python, die numerische Programmierung kognitiv n\"aher an die mathematische
Notation bringt. Durch den Einsatz eines Pr\"aprozessors, der vier klar definierte
syntaktische Muster transformiert, wird die extrinsische kognitive Last bei der
Implementierung von Matrix-Algorithmen auf ein Minimum reduziert.

Die Kernbeitr\"age des Projekts sind:

\begin{itemize}
	\item Die Klassen \texttt{Array2D} und \texttt{Array1D} mit sauberem Call-Protokoll
	\item Der Pr\"aprozessor, der \texttt{a(i,j) = expr} in \texttt{a.\_\_setcall\_\_(i,j,expr)}
	      \"ubersetzt
	\item Die korrekte Behandlung von Verbundzuweisungen und Lesezugriffen
	\item Das Leitbeispiel der Crout-LU-Zerlegung als Nachweis der praktischen Anwendbarkeit
\end{itemize}

\subsection{Kognitionspsychologische Schlussfolgerung}

Die psychologische Rechtfertigung des Projekts st\"utzt sich auf sechs unabh\"angige
Theorien:

\begin{itemize}
	\item \textbf{Arbeitsged\"achtnis (Miller):} Reduzierung der Chunk-Zahl pro Ausdruck
	\item \textbf{Kognitive Last (Sweller):} Elimination extrinsischer Last durch Notation
	\item \textbf{Mental Models (Johnson-Laird):} Direkte Abbildung des mathematischen
	      Modells in Code
	\item \textbf{Repr\"asentationskongruenz:} Strukturgleichheit zwischen Notation
	      und mentalem Modell
	\item \textbf{Perzeptuelle Transparenz (Blackwell):} Visuelle Struktur spiegelt
	      Semantik direkt wider
	\item \textbf{Flow-Theorie (Csikszentmihalyi):} Beseitigung von Unterbrechungsquellen
	      f\"ur den Implementierungs-Flow
\end{itemize}

\subsection{Ausblick}

M\"ogliche Erweiterungen des Projekts umfassen:

\begin{itemize}
	\item Verallgemeinerung auf $n$-dimensionale Arrays
	\item Jupyter-/IPython-Integration f\"ur interaktiven Einsatz
	\item Direktcompilierung nach C oder Fortran f\"ur Hochperformanz-Anwendungen
	\item Typsystem-Plugin f\"ur statische Analyse
	\item Erweiterung der Tupel-Konversion auf weitere Sprachkonstrukte (z.\ B.\
	      Comprehensions in Klammer-Syntax)
	\item Parallelisierung gro\ss{}er Array-Operationen durch automatische
	      \texttt{multiprocessing}-Annotation
\end{itemize}

\subsection{Implementierungshinweis}

Der Code ist verf\"ugbar unter:
\url{https://gitlab.com/epp-group/python}

\begin{thebibliography}{99}

\bibitem{miller1956}
G.~A. Miller.
\newblock The magical number seven, plus or minus two: Some limits on our
  capacity for processing information.
\newblock \emph{Psychological Review}, 63(2):81--97, 1956.

\bibitem{sweller1988}
J.~Sweller.
\newblock Cognitive load during problem solving: Effects on learning.
\newblock \emph{Cognitive Science}, 12(2):257--285, 1988.

\bibitem{sweller1994}
J.~Sweller, P.~Chandler, P.~Tierney, and M.~Cooper.
\newblock Cognitive load as a factor in the structuring of technical material.
\newblock \emph{Journal of Experimental Psychology: General}, 119(2):176--192,
  1990.

\bibitem{johnson-laird1983}
P.~N. Johnson-Laird.
\newblock \emph{Mental Models: Towards a Cognitive Science of Language,
  Inference, and Consciousness}.
\newblock Cambridge University Press, Cambridge, 1983.

\bibitem{festinger1957}
L.~Festinger.
\newblock \emph{A Theory of Cognitive Dissonance}.
\newblock Stanford University Press, Stanford, 1957.

\bibitem{csikszentmihalyi1990}
M.~Csikszentmihalyi.
\newblock \emph{Flow: The Psychology of Optimal Experience}.
\newblock Harper and Row, New York, 1990.

\bibitem{blackwell2001}
A.~F. Blackwell.
\newblock Pictorial representation and metaphor in visual language design.
\newblock \emph{Journal of Visual Languages and Computing}, 12(3):223--252,
  2001.

\bibitem{teitelman1966}
W.~Teitelman.
\newblock \emph{PILOT: A Step toward Man-Computer Symbiosis}.
\newblock PhD Thesis, MIT, Cambridge, 1966.

\bibitem{green1989}
T.~R.~G. Green.
\newblock Cognitive dimensions of notations.
\newblock In A.~Sutcliffe and L.~Macaulay, editors, \emph{People and Computers
  V}, pages 443--460. Cambridge University Press, 1989.

\bibitem{pane2001}
J.~F. Pane and B.~A. Myers.
\newblock Usability issues in the design of novice programming systems.
\newblock Technical Report CMU-CS-96-132, Carnegie Mellon University, 1996.

\bibitem{burnett2001}
M.~Burnett and C.~Baker.
\newblock A classification system for visual programming languages.
\newblock \emph{Journal of Visual Languages and Computing}, 5(3):287--300,
  1994.

\bibitem{dijkstra1968}
E.~W. Dijkstra.
\newblock Go to statement considered harmful.
\newblock \emph{Communications of the ACM}, 11(3):147--148, 1968.

\end{thebibliography}

\end{document}
